\documentclass[10pt,a4paper]{amsart}
\usepackage[margin=19mm]{geometry}
\usepackage{amsmath,amssymb,mathtools}
\usepackage[scr=rsfs,cal=euler]{mathalfa}
\usepackage{xfrac}
\usepackage[unicode,hidelinks,bookmarks=false]{hyperref}
\newcommand{\ie}{i.e.\ }
\newcommand{\cf}{cf.\ }
\newcommand{\resp}{respectively\ }
\newcommand{\Subsection}{\subsection}
\providecommand{\nobreakdash}{\nobreak\hbox{-}}
\setlength{\emergencystretch}{2em}
\multlinegap0pt
\usepackage{amssymb}
\usepackage{mathtools}
\usepackage{bm}
\usepackage[shortlabels]{enumitem}
\usepackage{algorithm}\usepackage{algpseudocode}
\usepackage{xcolor}
\usepackage{float}
\newtheorem{proposition}{Proposition}
\newcommand{\calC}{\mathcal C}
\newcommand{\hR}{\widehat R}
\newcommand{\hU}{\widehat U}
\newcommand{\hphi}{\widehat\phi}
\newcommand{\norm}[1]{\lVert #1\rVert}
\newcommand{\tU}{\widetilde U}
\title{The Resolution of Causal Heterogeneity}
\author{Yuki Ohnishi$^{*}$ and Fan Li$^{\dagger}$}
\date{Source: arXiv:2607.17280v1 (CC BY 4.0)}
\begin{document}
\maketitle

\noindent\emph{Source and licence.} The Resolution of Causal Heterogeneity. Version arXiv:2607.17280v1 (CC BY 4.0), distributed by arXiv under the Creative Commons Attribution 4.0 International licence, http://creativecommons.org/licenses/by/4.0/, as recorded in the arXiv licence metadata for this exact version and shown on the abstract page https://arxiv.org/abs/2607.17280v1. Full licence notice: \texttt{licenses/arxiv-2607.17280-CC-BY-4.0.txt}.\par\smallskip

\noindent\textit{Editorial scope.} Counted unit: the article's proposition prop:pseudo (source0.tex lines 952--972). The article's complete original proof of it is delivered with it (source0.tex lines 974--997, one \texttt{proof} environment of 24 lines). No mathematics is written, repaired or re-derived: the statement and the proof are the article's own lines, in the article's own notation, and no retained line is substituted or re-flowed.  No further source-local context is needed: every cross-reference of every retained line resolves inside the two delivered spans.  Nothing was omitted from any retained span, so the package carries no omission record.  No retained line contains a citation, so the wrapper carries no bibliography and no bibliography entry.  No retained line contains a figure environment, an image-inclusion command or drawing code, so no image is delivered and the package declares no asset dependency.  Editorial formatting only, in addition: an AMS \texttt{amsart} preamble that restates the article's own declarations and macros used by the retained lines, namely the environments \texttt{proposition} on separate counters, together with the standard packages those lines need.  The retained environments print in this document as Proposition 1 in order of appearance, whereas the article prints the same items with the numbering of its own full document.  The complete native source of the article as distributed in the arXiv e-print for this exact version is bundled unchanged as \texttt{sources/205552/source0.tex}.
\begin{proposition}[Pseudo-feature identity]\label{prop:pseudo}
For every codebook $c\in\calC^K$ and every observation $i$,
\begin{equation}\label{eq:pseudoidentity}
\hphi_{g_c,i}
=
\bigl\|\tU_i-c_{h_c(\hU_i)}\bigr\|^2
-
\bigl\|H\hR_i\bigr\|^2,
\end{equation}
where $h_c(\hU_i)$ is the Voronoi cell of $\hU_i$, not of $\tU_i$.  Consequently, for any nonnegative weights $w=(w_1,\dots,w_n)$,
\begin{equation}\label{eq:pseudoweighted}
\frac1n\sum_{i=1}^n w_i\hphi_{g_c,i}
=
\frac1n\sum_{i=1}^n w_i
\bigl\|\tU_i-c_{h_c(\hU_i)}\bigr\|^2
-
\frac1n\sum_{i=1}^n w_i
\bigl\|H\hR_i\bigr\|^2 .
\end{equation}
The second term is free of $c$ and therefore can be dropped for optimization over codebooks, but it must be retained when reporting corrected risk values.
\end{proposition}

\begin{proof}
For the quantization loss $g_c(u)=\min_{h\le K}\norm{u-c_h}^2$, off Voronoi boundaries,
\[
\nabla g_c(u)=2\{u-c_{h_c(u)}\}.
\]
Thus the corrected evaluation at observation $i$ is
\[
\hphi_{g_c,i}
=
\norm{\hU_i-c_{h_c(\hU_i)}}^2
+
2\{\hU_i-c_{h_c(\hU_i)}\}^\top H\hR_i .
\]
Apply the identity
\[
\norm{a-c}^2+2(a-c)^\top v
=
\norm{a+v-c}^2-\norm{v}^2
\]
with $a=\hU_i$
, $v=H\hR_i$
, and $c=c_{h_c (\hU_i)}$.
This proves \eqref{eq:pseudoidentity}. Multiplying by weights and summing gives \eqref{eq:pseudoweighted}.
\end{proof}

\end{document}
